\ifdefined\SUComplete\else
\documentclass[UTF8,11pt,fontset=none]{ctexart}
\usepackage[a4paper,margin=25mm]{geometry}
\setCJKmainfont{Droid Sans Fallback}[AutoFakeBold=2]
\setmainfont{DejaVu Serif}
\usepackage{amsmath,amssymb,mathtools,booktabs,longtable,xcolor,hyperref}
\hypersetup{unicode,colorlinks=true,linkcolor=blue!45!black,urlcolor=blue!45!black}
\allowdisplaybreaks[2]
\setlength{\emergencystretch}{2em}
\newcommand{\SU}{\mathrm{SU}(2)}
\newcommand{\Dic}{\mathrm{Dic}}
\newcommand{\Tr}{\operatorname{Tr}}
\newcommand{\Vol}{\operatorname{Vol}}
\newcommand{\Arg}{\operatorname{Arg}}
\newcommand{\Li}{\operatorname{Li}}
\title{SU(2) 反常与 ADE 有限子群：从有向体积到具体闭余循环}
\author{}
\date{2026 年 9 月 22 日}

\begin{document}
\maketitle
\begin{abstract}
本文是对 SU(2) 背景反常 ADE 限制问题的审计和补充计算。
首先固定源笔记 A 第 10 页式 (7.1) 的相位方向，说明它与仓库原有正体积代表是逆余循环。
随后证明全局有限填充给出 closed 三余循环，并用 $C_4$ 的 bar 三循环证明 $k=1$ 非恰当。
对二元二面体群给出精确群律、八分支有向体积表和有理数闭性验证；
从三个群元素的六个迹直接给出边长体积公式；
对 $2T,2O,2I$ 写出精确元素集合、有限分支相位公式和三组独立手算实例。
另给有界的精确代数计数代表、与几何代表的显式余边界及完整类阶证明。
数值结果区分数学证明、精确代数检查和高精度浮点检查。
\end{abstract}

\tableofcontents
\clearpage
\fi

\section{资料审计与统一约定}

沿用仓库记录的源 A 笔记式 (7.1) 约定（原附件当前不在 checkout）：
\begin{equation}
 \omega_k^{A}(g_1,g_2,g_3)=
 \exp\left(-\frac{ik}{\pi}\Vol_A(C)\right),
 \qquad \Vol_A(C)=\int_C dV_{S^3}.
 \label{eq:source-sign}
\end{equation}
同一笔记式 (1.1)、(2.12) 使用
\[
 \omega_3^A=-\frac{dV_{S^3}}{2\pi^2},
 \qquad \int_{S^3}\omega_3^A=-1.
\]
仓库原有接口选取 $+dV_{S^3}/(2\pi^2)$ 和 $\exp(+ikV/\pi)$。
因此本报告的新接口满足
\begin{equation}
 \omega_k^{A}=\left(\omega_k^{\rm old}\right)^{-1},
 \label{eq:inverse-convention}
\end{equation}
这是初末网络比值的代表转换，不是 level 位移，也不是未经说明地改变群乘法。

仍采用
\[
g=x_0I+i\sum_{a=1}^3x_a\sigma_a,\qquad
F=dA+A\wedge A,\qquad
I_4=-\frac{k}{8\pi^2}\Tr(F\wedge F).
\]
群元累计顶点为
\[
p_0=e,\quad p_1=g_1,\quad p_2=g_1g_2,\quad p_3=g_1g_2g_3.
\]
\section{三个群元素直接给出六边长及体积}

这是人类 review 要求的计算路线：从三个群元素出发，经六条边长直接得到体积。
设 $q(p)$ 为群元的四元数坐标，因
$q(p)\cdot q(r)=\tfrac12\Tr(p^{-1}r)$，六条短测地线边长按用户边序 $(01,02,03,23,13,12)$ 为
\begin{equation}
\boxed{\begin{aligned}
\ell_1=\ell_{01}&=\arccos\tfrac12\Tr g_1,&
\ell_2=\ell_{02}&=\arccos\tfrac12\Tr(g_1g_2),\\
\ell_3=\ell_{03}&=\arccos\tfrac12\Tr(g_1g_2g_3),&
\ell_4=\ell_{23}&=\arccos\tfrac12\Tr g_3,\\
\ell_5=\ell_{13}&=\arccos\tfrac12\Tr(g_2g_3),&
\ell_6=\ell_{12}&=\arccos\tfrac12\Tr g_2 .
\end{aligned}}\label{eq:group-six-lengths}
\end{equation}
对边为 $j$ 与 $j+3$。这些是单位 $S^3$ 上的弧长，不能代入欧氏 Cayley--Menger 体积。

这六条边长的完整体积函数，是第 \ref{sec:symbolic-ade} 节的式 \eqref{eq:exact-master-phase}。
它是 Murakami 定理 1.2\cite{murakami} 的群元代入形式，固定 $z$ 的全部偏导都已展开，
取向由 $D=\tfrac14\Tr[A(BC-CB)]$（$A=g_1$，$B=g_1g_2$，$C=g_1g_2g_3$）的符号给出。
当三个群元属于 ADE 有限群时，第 \ref{sec:polar-dual} 节的式 \eqref{eq:polar-master}
进一步给出不含 $\Li_2$ 的精确值。
六条边长不能区分镜像，$D=0$ 时也不能一概置零；完整分支见第 \ref{sec:E-explicit} 节。

\paragraph{与历史接口的边序换算。}
历史数值接口 \texttt{edge\_lengths\_quaternions} 与本报告后半部保留的原始推导使用旧边序
$(23,13,03,01,02,12)$，它依次对应上面的 $(\ell_4,\ell_5,\ell_3,\ell_1,\ell_2,\ell_6)$。
全文新结果一律使用用户边序；旧边序只出现在历史接口与原推导中。


\section{closed 与 non-exact}

单位顶点规范化光滑奇异链上的边、面和三维链满足
\begin{align}
\partial E(a,b)&=[b]-[a],\\
\partial F(a,b,c)&=E(b,c)-E(a,c)+E(a,b),\\
\partial T(a,b,c,d)&=F(b,c,d)-F(a,c,d)+F(a,b,d)-F(a,b,c).
\end{align}
这里的 E/F/T 使用相邻重复退化、精确凸包判定和有限锥点规则；一般降秩链保留。
对五个累计顶点组成的五项三链 $Z$，面逐项相消，故 $\partial Z=0$。
由于 $\int_{S^3}dV=2\pi^2$，有
\[
\int_ZdV\in2\pi^2\mathbb Z,
\qquad \delta\omega_k^A=1.
\]
这一步是全局证明，不依赖随机五边形测试。

取 $h=e^{2\pi i\sigma_1/n}$，定义进位函数
\[
N_n(a,b)=\left\lfloor\frac{a+b}{n}\right\rfloor,
\qquad [a+b]_n=a+b-nN_n(a,b).
\]
规范化 bar 三循环为
\[
z_n=\sum_{r=0}^{n-1}[h\mid h^r\mid h].
\]
A 半球约定的限制代表是
\begin{equation}
\omega_{k,C_n}(h^a,h^b,h^c)=
\exp\left(\frac{2\pi i k a}{n}N_n(b,c)\right).
\label{eq:Cn-cocycle}
\end{equation}
因此
\begin{equation}
P_n(\omega_k)=\prod_{r=0}^{n-1}\omega_k(h,h^r,h)=e^{2\pi ik/n}.
\label{eq:Cn-pairing}
\end{equation}
若 $\omega_1=\delta\beta$，限制到 $C_4$ 后仍为余边界，而余边界在 $z_4$ 上配对为 1；式 (\ref{eq:Cn-pairing}) 却给出 $P_4=i$。所以
\[
\boxed{\omega_1\text{ closed but not exact}.}
\]
这里应区分半球代表与几何 E/F/T 代表的逐点值。比较两套边链时，先取
$\partial A(g)=\gamma_A(g)-\gamma_{\rm EFT}(g)$，再取
\[
\partial B(g,h)=S_A(g,h)-S_{\rm EFT}(g,h)
-\bigl(gA(h)-A(gh)+A(g)\bigr).
\]
$H_1(S^3)=H_2(S^3)=0$ 保证这些比较链存在；
相位之比为 $\delta\beta$，其中
\[
\beta(g,h)=\exp\left[-\frac{ik}{\pi}\int_{B(g,h)}dV\right].
\]
因此三循环配对相同，不要求代表逐点相等。

\section{二元二面体群的具体相位}

令
\[
U=e^{i\pi\sigma_1/n},\qquad V=i\sigma_2,
\qquad U^{2n}=1,\quad V^2=U^n,\quad VUV^{-1}=U^{-1}.
\]
使用标签 $V^sU^r$，其中 $r\in\{0,\ldots,2n-1\}$、$s\in\{0,1\}$。
旧仓库内部标签写成 $U^rV^s$ 时，需在 $s=1$ 分支将 $r$ 换成 $-r$。

定义 $[x]_{2n}\in\{0,\ldots,2n-1\}$ 和
\[
N_n(a,b)=\left\lfloor\frac{[a]_{2n}+[b]_{2n}}{2n}\right\rfloor.
\]
源笔记 A 的半球构造给出 $\Vol_A=\pi^2v_n$，其中
\begin{longtable}{@{}c p{0.65\linewidth}@{}}
\toprule
$(s_1s_2s_3)$ & $v_n$\\ \midrule
000 & $-r_1N_n(r_2,r_3)/n$\\
001 & $([r_3-r_1-r_2]_{2n}-n)N_n(r_1,r_2)/n$\\
010 & $r_1r_3/(2n^2)$\\
011 & $-([r_2-r_1]_{2n}-n)N_n(r_1,n-r_2+r_3)/n$\\
100 & $-r_1N_n(r_2,r_3)/n$\\
101 & $[n-r_1-r_2+r_3]_{2n}r_2/(2n^2)$\\
110 & $(r_1-n)N_n(n-r_1+r_2,r_3)/n$\\
111 & $-[n-r_1+r_2]_{2n}[n-r_2+r_3]_{2n}/(2n^2)$\\
\bottomrule
\end{longtable}
最终相位为
\[
\boxed{\omega_{k,\Dic_n}(V^{s_1}U^{r_1},V^{s_2}U^{r_2},V^{s_3}U^{r_3})
=e^{-i\pi k v_n}.}
\]
这是一套完整三元组规则，不是只列代表性元素。

程序 \texttt{scripts/dic\_formula.py} 实现表格，\texttt{scripts/check\_dic\_table.py} 用有理数检查
\[
\delta v_n\in2\mathbb Z.
\]
对于 $n=2,3,4,5$，分别穷举 $4096,20736,65536,160000$ 个四元组，全部通过。
有限穷举不构成任意 $n$ 的符号证明。旧交接只概述使用进位恒等式约去，
未展示逐分支消去或此表对应的完整边面链；因此本报告将这一证明列为待补。
将该表识别为几何 E/F/T 的同类代表，还需落实它所用的几何链，
再应用上节的边、面比较；不能只从两者都闭就推出同类。

\section{E 型群元素的完整几何相位公式}\label{sec:E-explicit}

\subsection{三个精确元素集合}
令 $e_0,\ldots,e_3$ 为四维坐标基，$\varphi=(1+\sqrt5)/2$。
用下面的集合直接表示 $E_6,E_7,E_8$ 所对应的二元群：
\begin{align*}
\mathcal Q_6={}&\{\pm e_r:0\le r\le3\}
\cup\{(\epsilon_0,\epsilon_1,\epsilon_2,\epsilon_3)/2:\epsilon_r=\pm1\},\\
\mathcal Q_7={}&\mathcal Q_6\cup
\{(\epsilon e_r+\epsilon'e_s)/\sqrt2:r<s,\ \epsilon,\epsilon'=\pm1\},\\
\mathcal Q_8={}&\mathcal Q_6\cup
\{\rho(0,\epsilon_1/2,\epsilon_2\varphi/2,\epsilon_3\varphi^{-1}/2):
\rho\in A_4,\ \epsilon_r=\pm1\}.
\end{align*}
其中 $A_4$ 表示偶坐标置换。三群阶分别为 $24,48,120$。
在 $g=q_0I+i\mathbf q\cdot\boldsymbol\sigma$ 约定下，乘法是
\[
(a,\mathbf u)(b,\mathbf v)
=(ab-\mathbf u\cdot\mathbf v,\ a\mathbf v+b\mathbf u-\mathbf u\times\mathbf v).
\]
所有乘积、相等、秩、正负号判定均在 $\mathbb Q$、
$\mathbb Q(\sqrt2)$ 或 $\mathbb Q(\sqrt5)$ 中完成，平方根取正根。

\subsection{所有三元组的有限分支}
以下定义相位公式中每一求和项。$[v_0,\ldots,v_j]$ 是单位化顶点的径向有向单形；
计算时可保存未归一化的正射线。令
\[
S(v_0,\ldots,v_j)\iff0\notin\operatorname{conv}\{v_0,\ldots,v_j\}.
\]
等价地，不存在非零全非负系数满足 $\sum_i\lambda_iv_i=0$。
下列各链遇相邻相同顶点均定义为零；其余情形固定 $J=e_1$，取
\[
E(a,b)=\begin{cases}[a,b],&S(a,b),\\
[a,aJ]+[aJ,b],&\text{否则}.
\end{cases}
\]
对 $Z=\sum_{\nu=1}^M c_\nu[v_{\nu0},\ldots,v_{\nu r}]$ 定义
\begin{align*}
t_*&=\min\{t\in\{0,\ldots,3M\}:a(1,t,t^2,t^3)
\notin\operatorname{span}(v_{\nu0},\ldots,v_{\nu r})\ \forall\nu\},\\
C_a(Z)&=\sum_\nu c_\nu[a(1,t_*,t_*^2,t_*^3),v_{\nu0},\ldots,v_{\nu r}],
\qquad C_a(0)=0.
\end{align*}
于是
\[
F(a,b,c)=\begin{cases}[a,b,c],&S(a,b,c),\\
C_a(E(b,c)-E(a,c)+E(a,b)),&\text{否则},\end{cases}
\]
\[
T(a,b,c,d)=\begin{cases}[a,b,c,d],&S(a,b,c,d),\\
C_a(F(b,c,d)-F(a,c,d)+F(a,b,d)-F(a,b,c)),&\text{否则}.
\end{cases}
\]
每个待避真子空间在三次矩曲线上至多排除三个整数，故 $F$ 至多 6 项、
19 个候选，$T$ 至多 24 项、73 个候选。锥顶点在原单形张成之外，
因而安全单形锥起后仍安全。降秩项保留在链中，其三维积分为零。

给定任意 $q_1,q_2,q_3\in\mathcal Q_m$，展开
\[
T(e_0,q_1,q_1q_2,q_1q_2q_3)
=\sum_{\nu=1}^{M\le24}c_\nu[v_{\nu0},v_{\nu1},v_{\nu2},v_{\nu3}].
\]
记 $D_\nu=\det(v_{\nu0},\ldots,v_{\nu3})$，每项边长为
\[
\ell_{\nu,ij}=\arccos\frac{v_{\nu i}\cdot v_{\nu j}}
{\|v_{\nu i}\|\|v_{\nu j}\|},\qquad
(ij)=(01,02,03,23,13,12).
\]
由式 \eqref{eq:exact-master-phase} 中已全部展开的体积函数 $W$ 得到
\begin{equation}
\boxed{\omega_{E_m,k}(q_1,q_2,q_3)=
\exp\!\left[-\frac{ik}{\pi}\sum_{\nu:D_\nu\ne0}
c_\nu\operatorname{sgn}(D_\nu)W(\ell_{\nu1},\ldots,\ell_{\nu6})\right],
\quad m=6,7,8.}
\end{equation}
所有符号和求和项均已由群元定义。非退化输入的专属短表见第 \ref{sec:polar-dual} 节和第 \ref{sec:complete-gram-table} 节；退化输入的初等形式见式 \eqref{eq:degenerate-elementary}。
左乘保持取向、秩与锥点选择，配合 $\partial T=\delta F$，给出全局五边形。

\subsection{三组有独立手算推导的例子}
共同取 $g_1=e_1=i\sigma_1$、$g_2=e_3=i\sigma_3$，则 $p_2=e_2$。
\begin{align*}
2T:\quad g_3&=(1,1,-1,-1)/2,& \omega_{A,k}&=e^{-ik\pi/32},\\
2O:\quad g_3&=(1,1,0,0)/\sqrt2,& \omega_{A,k}&=e^{-ik\pi/16},\\
2I:\quad g_3&=(\varphi/2,\varphi^{-1}/2,0,-1/2),&
\omega_{A,k}&=\exp\!\left[-\frac{ik}{2}\arctan\frac1{4\varphi+1}\right].
\end{align*}
第一行的 $p_3=(1,1,1,1)/2$ 将正交球面四面体（体积 $\pi^2/8$）
分成四个全等部分；第二行的 $p_3=(0,0,1,1)/\sqrt2$ 将它平分。
第三行的 $p_3=(0,1/2,\varphi/2,\varphi^{-1}/2)$ 与 $e_1,e_2$ 构成 $S^2$ 三角形，面积
\[
\Omega=2\arctan\frac{\varphi^{-1}}{3+\varphi}
=2\arctan\frac1{4\varphi+1}.
\]
它与 $e_0$ 的球面连接体积为
$V=\Omega\int_0^{\pi/2}\sin^2u\,du=\pi\Omega/4$，得上述相位。
第二、三行确实使用超出共享 $2T$ 子群的元素。

\subsection{类阶与点值的区别}
几何代表与基本 CS/String 类的识别使用归一化三形式及自然性。
第 \ref{sec:E-algebraic} 节进一步给出 Gysin 序列和障碍类的直接证明，
与 Epa--Ganter 定理 1.1\cite{epa} 一致：
\[
\operatorname{ord}[\omega_k|_\Gamma]=\frac{|\Gamma|}{\gcd(k,|\Gamma|)},
\qquad [\omega_k|_\Gamma]=0\iff |\Gamma|\mid k.
\]
E 型的类周期分别为 $24,48,120$；这不要求上述几何代表逐点为对应阶的单位根，
也不是用几个数值相位证明类阶。

\section{历史比较：E 型代数代表与显式余边界}\label{sec:E-algebraic}

\noindent\textbf{09-23 路线说明：}本节保留旧研究过程，不作为本轮用户要求的解法。
当前主结果见第 \ref{sec:symbolic-ade}--\ref{sec:oi-exact} 节，保持原几何相位，
不采用 transfer，也不以计数代表或数值比较替代实际体积。

本节保留第 \ref{sec:E-explicit} 节的几何代表，另给同类的计数代表。
后者只使用原字段中的四阶行列式、精确符号和有界整数求和，
不含体积积分、反三角函数或二重对数。它并非原几何相位的逐点化简，
也不是已经导出的 E 型专属短分支表。

\subsection{由三个群元素到整数指数的完整公式}
取 $\Gamma=\mathcal Q_6,\mathcal Q_7,\mathcal Q_8$，分别令 $N=24,48,120$。
使用前节全部给定的群元素、乘法、相邻重复规则及 E/F/T 有限链，展开
\[
T(e_0,g_1,g_1g_2,g_1g_2g_3)
=\sum_{\nu=1}^{M\le24}c_\nu[v_{\nu0},v_{\nu1},v_{\nu2},v_{\nu3}].
\]
记 $D_\nu=\det(v_{\nu0},v_{\nu1},v_{\nu2},v_{\nu3})$。
对 $D_\nu\ne0$、$h\in\Gamma$ 和 $0\le j,s\le3$，定义 Cramer 分子
\[
d_{\nu hjs}=\det(v_{\nu0},\ldots,v_{\nu,j-1},he_s,
v_{\nu,j+1},\ldots,v_{\nu3}).
\]
令 $\operatorname{lsign}(a_0,a_1,a_2,a_3)$ 为首个非零项的符号，并设
\begin{equation}
I_{\nu h}=\prod_{j=0}^3\mathbf1\!\left\{
\operatorname{lsign}(d_{\nu hj0},d_{\nu hj1},d_{\nu hj2},d_{\nu hj3})
=\operatorname{sgn}D_\nu\right\}.
\label{eq:count-indicator}
\end{equation}
各系数组不会全为零，因为非零 Cramer 线性泛函不能湮灭基
$he_0,\ldots,he_3$。完整相位为
\begin{equation}
A_\Gamma(g_1,g_2,g_3)
=\sum_{\nu:D_\nu\ne0}c_\nu\operatorname{sgn}(D_\nu)
\sum_{h\in\Gamma}I_{\nu h}\in\mathbb Z.\label{eq:count-integer}
\end{equation}
\begin{equation}
\boxed{\widehat\omega_{\Gamma,k}(g_1,g_2,g_3)
=\exp\left(-\frac{2\pi i k}{N}A_\Gamma(g_1,g_2,g_3)\right),
\quad k\in\mathbb Z.}\label{eq:count-phase}
\end{equation}
实现可直接返回 $-kA_\Gamma\bmod N$，复数近似不是定义的一部分。
令 $V_\nu$ 以四个顶点为列，$L_h$ 为左乘 $h$ 的矩阵；等价判据为
$V_\nu^{-1}L_h$ 的每一行首个非零项均为正。
每次至多 $24N$ 个单形与轨道点的判定，即 $576,1152,2880$ 个判定；
此数不是基本算术运算次数。

全部算术留在 $\mathbb Q$、$\mathbb Q(\sqrt2)$、$\mathbb Q(\sqrt5)$ 中。
对 $a+b\sqrt d$，同号时直接取符号，异号时比较 $a^2$ 与 $db^2$；
$d=2,5$ 非平方，非零有理 $a,b$ 不会造成相等歧义。
正射线缩放不改变结果，负比例不能视为相同顶点。
降秩项在链级保留，仅在三维计数与积分中贡献零；
也不能在参数化奇异链中把倒序单形直接当成负单形约去。

\subsection{全部输入的闭性、归一化与退化分支}
令 $r(\varepsilon)=(1,\varepsilon,\varepsilon^2,\varepsilon^3)$。
Cramer 法则给出
\[
\bigl(V_\nu^{-1}L_h r(\varepsilon)\bigr)_j
=D_\nu^{-1}\sum_{s=0}^3d_{\nu hjs}\varepsilon^s.
\]
充分小的正 $\varepsilon$ 下，其符号就是式 \eqref{eq:count-indicator}。
群和全部输入有限，涉及的边界、降秩像及其群平移只占有限个真线性子空间。
任一非零线性泛函作用在 $r(\varepsilon)$ 上都是次数至多三的非零多项式。
故存在一个共同的充分小正实数，使全部符号规则同时成立且避开这些低维像。
这提供普通实球面中的解释；实际公式无需选择浮点 $\varepsilon$。

满秩安全径向单形在单位点 $y$ 上的局部重数为
$\operatorname{sgn}D_\nu$ 当且仅当 $V_\nu^{-1}y$ 的四坐标全正。
因此 $A_\Gamma$ 是 $T$ 在完整轨道
$\{hr(\varepsilon)/\|r(\varepsilon)\|:h\in\Gamma\}$ 上的有向重数之和。
左乘保向并置换轨道，故此计数的齐次版本等变。
对闭三链 $Z$，在避开边界和低维像的点上，局部次数定理给出
\begin{equation}
\chi_y(Z)=\deg Z=\frac1{2\pi^2}\int_ZdV\in\mathbb Z.
\label{eq:count-degree}
\end{equation}
规范化奇异链的分裂提升只添加退化项，低维像对该重数与三维积分都为零，
所以同一公式适用于这里的规范化链。

由 $\partial T=\delta F$，有 $\partial(\delta T)=0$，从而
\begin{equation}
\boxed{\delta A_\Gamma=N\deg(\delta T)\in N\mathbb Z.}
\label{eq:count-cocycle}
\end{equation}
非齐次写成
\[
A(b,c,d)-A(ab,c,d)+A(a,bc,d)-A(a,b,cd)+A(a,b,c)\equiv0\pmod N.
\]
式 \eqref{eq:count-phase} 因而满足完整五边形。
任一输入为单位元时累计顶点相邻重复，$T=0$、$A=0$、相位为 1。
反足及无反足但正依赖的输入由同一 E/F/T 分支覆盖。

\subsection{与几何代表逐点成立的转换公式}
定义实上链
\[
t(g,h,l)=\frac1{2\pi^2}\int_{T(e_0,g,gh,ghl)}dV,
\qquad \tau(g,h,l)=\frac{A_\Gamma(g,h,l)}N,
\qquad f=\tau-t.
\]
式 \eqref{eq:count-degree} 给出 $\delta\tau=\delta t=\deg(\delta T)$，
故 $\delta f=0$ 是实数等式。明确取有限和
\begin{equation}
B(g,h)=\frac1N\sum_{x\in\Gamma}f(x,g,h),
\qquad \beta_k(g,h)=\exp[-2\pi i kB(g,h)].
\label{eq:count-beta}
\end{equation}
对 $\delta f(x,g,h,l)=0$ 求平均，使用 $x\mapsto xg$ 置换群，得到
\[
f(g,h,l)=B(h,l)-B(gh,l)+B(g,hl)-B(g,h)=\delta B(g,h,l).
\]
因此不是仅断言存在某个二余链，而是逐点有
\begin{equation}
\boxed{\widehat\omega_{\Gamma,k}=\omega^{\rm geom}_{A,k}\,\delta\beta_k,
\qquad \delta\beta_k(g,h,l)=
\frac{\beta_k(h,l)\beta_k(g,hl)}{\beta_k(gh,l)\beta_k(g,h)}.}
\label{eq:count-conversion}
\end{equation}
$B$、$\beta$ 归一化。计算新相位本身不需先计算 $B$；
只有比较两代表时，式 \eqref{eq:count-beta} 才调用已给出的几何体积公式。

\subsection{完整类阶的 Gysin 证明}
$\Gamma$ 在 $S^3$ 上左乘自由且保向。相应定向四维实表示的球面丛满足
\[
S^3\longrightarrow E\Gamma\times_\Gamma S^3\longrightarrow B\Gamma,
\qquad E\Gamma\times_\Gamma S^3\simeq S^3/\Gamma.
\]
商为定向闭三流形，Gysin 序列中的相关部分是
\[
H^3(S^3/\Gamma;\mathbb Z)\cong\mathbb Z
\xrightarrow{\ \times N\ }H^0(B\Gamma;\mathbb Z)\cong\mathbb Z
\xrightarrow{\ \smile e\ }H^4(B\Gamma;\mathbb Z)\longrightarrow0.
\]
第一箭头是沿纤维积分：在上述同伦等价下纤维包含对应覆盖
$S^3\to S^3/\Gamma$，其次数为 $N$；最右零来自商空间的四次上同调消失。
所以 $H^4(B\Gamma;\mathbb Z)=\mathbb Z/N$，由 Euler 类 $e$ 生成。

利用 $S^3$ 的二连通性可在 bar 模型三骨架上构造等变参考截面；
其四胞腔延拓障碍是边界映到球面的次数，即 Euler 类。
E/F/T 的积分链与该参考截面的链在零至二维依次比较，
差由更高一维链填充；三维差的次数是一个整数三上链。
因此两种四维障碍只差这个三上链的余边界，
$[\deg(\delta T)]=e$（随所选定向固定符号）。
这一比较不把形式链假定为一个单独的胞腔映射。
由 $0\to\mathbb Z\to\mathbb R\to U(1)\to0$ 的连接同态，
$[\exp(-2\pi ikt)]$ 映到 $-k[\deg(\delta T)]$。
有限群的正次实上同调为零，此连接同态为同构。结合转换公式，得到
\begin{equation}
\boxed{\operatorname{ord}[\widehat\omega_{\Gamma,k}]
=\operatorname{ord}[\omega^{\rm geom}_{A,k}]
=\frac{N}{\gcd(k,N)}.}\label{eq:count-order}
\end{equation}
这也与 Epa--Ganter 的结果一致\cite{epa}。
$C_4$ 配对只能校准该限制；三群循环子群阶的最小公倍数为 $12,24,60$，
故循环限制本身不能证明完整的 $24,48,120$ 阶。

\subsection{计数为 1、2、4 的具体输入及代表差异}
仍取 $g_1=e_1,g_2=e_3$ 和前节的三个 $g_3$。
每条链均为正取向直接四面体，全部贡献者如下：
\begin{longtable}{@{}c p{0.60\linewidth} c@{}}
\toprule 群 & 满足 $I_{\nu h}=1$ 的全部 $h$ & $A$\\\midrule
$2T$ & $e_0$ & 1\\
$2O$ & $e_0,\ (1,0,1,0)/\sqrt2$ & 2\\
$2I$ & $e_0,\ (1,\varphi^{-1},\varphi,0)/2,$\newline
$(\varphi,1,\varphi^{-1},0)/2,\ (\varphi^{-1},\varphi,1,0)/2$ & 4\\
\bottomrule
\end{longtable}
相应计数相位为 $e^{-i\pi k/12},e^{-i\pi k/12},e^{-i\pi k/15}$，
而原几何相位为 $e^{-i\pi k/32},e^{-i\pi k/16},e^{-ik\theta/2}$，
其中 $\theta=\arctan(1/(4\varphi+1))$。
二者不同是式 \eqref{eq:count-conversion} 的代表变换，不是数值误差。

事实上最后一个原几何相位在非零整数 $k$ 时甚至不是单位根。
令 $u=(4\varphi+1)^{-1}$，则
\[
e^{2i\theta}=\frac{1+iu}{1-iu}\in K=\mathbb Q(\sqrt5,i),
\qquad 0<\theta<\pi/4.
\]
$K$ 中单位根只有 $\{\pm1,\pm i\}$：若含原始 $m$ 次根，则
$\phi(m)\mid4$，候选仅为 $1,2,3,4,5,6,8,10,12$。
三个二次子域 $\mathbb Q(\sqrt5),\mathbb Q(i),\mathbb Q(\sqrt{-5})$
排除需要 $\sqrt{-3},\sqrt2,\sqrt3$ 的候选；
$\mathbb Q(\zeta_5)$ 为循环四次扩张，而 $K$ 的 Galois 群为 $C_2\times C_2$，
故也排除 $5,10$。上述开区间排除剩余四根，故 $\theta/\pi$ 无理。
因此不能把原几何 E8 接口的点值改写为 $\mu_{120}$ 值表。


\section{复现结果与完成状态}\label{sec:ade-validation}

本轮新增的精确计数接口与验证记录为：
\begin{itemize}
\item 当前 39 项单元测试通过，覆盖原几何接口及新增计数、链边界、
正依赖、射线缩放、字段约化和输入验证；
\item $2T,2O,2I$ 各选取 16 个确定性四元组，全部精确验证
$\delta A\equiv0\pmod N$；保存五个整数计数、全部输入与链分支；
\item 三组手算例精确计数分别为 $1,2,4$；$C_4$ 的计数相位配对
指数分别为 $6,12,30$（分母为 $24,48,120$），均给出 $i$；
\item 每群一个完整的式 \eqref{eq:count-conversion} 比较，
保存四个 $\beta$ 的全部 $N$ 项平均输入。
\end{itemize}
转换比较使用 45 位输出精度，残差分别为
$4.5484781\times10^{-46}$、$3.1517170\times10^{-46}$、
$6.1152770\times10^{-46}$，均低于检查阈值 $10^{-39}$。
精确整数五边形没有浮点容差；上述转换比较仍是浮点验证，
代表转换和全输入闭性的依据是第 \ref{sec:E-algebraic} 节的证明。
完整数据见 \path{docs/review-0922/SU2_E_algebraic_validation.json}。
这些是选定样本，未穷举 E 型全部四元组。

以下浮点记录来自前一轮几何代表复核，作为历史基线保留；
不代表新增计数接口的当前测试总数或精确模整数覆盖数：
\begin{itemize}
\item 当时 26 项单元测试通过，包括六边长顺序、镜像取向、中心退化及三组手算例；
\item $2T,2O,2I$ 各 3 组五边形，保存精确输入和全部五项；
\item 各群的手算例与边长路线、独立二面角路线比较；
\item 中心 $(-1,-1,-1)$ 的相位为 $-1$，含单位元的相位为 1。
\end{itemize}
以 $k=1$、65 位输出记录计算，五项左右乘积的最大残差为
\begin{longtable}{@{}c c c@{}}
\toprule
群 & 四元组数 & 最大残差\\ \midrule
$2T$ & 3 & $1.143\times10^{-65}$\\
$2O$ & 3 & $9.108\times10^{-66}$\\
$2I$ & 3 & $8.464\times10^{-66}$\\
\bottomrule
\end{longtable}
完整数据保存在 \path{docs/review-0922/SU2_ADE_phase_validation.json}。
它逐组列出 $(a,b,c,d)$ 的精确四元数，以及
\[
\omega(b,c,d)\omega(a,bc,d)\omega(a,b,c)
=\omega(ab,c,d)\omega(a,b,cd)
\]
中的全部五个相位、各自累计点、分支、每个四面体的六边长及有向体积。
这不是 $24^4,48^4,120^4$ 个四元组的穷举，也不是区间算术误差证书。
另重跑 Q8 的全部 4096 个四元组，最大残差 $1.138\times10^{-86}$，
以及三群的精确乘法闭包；D 表有理数检查保留历史报告。
这些检查不算作 E 型四元组穷举。

\section{源码入口、修复与剩余边界}
\begin{sloppypar}
\begin{itemize}
\item \path{src/su2_omega.py}：\texttt{omega\_quaternions\_edge} 由三个群元经边长求 A 约定相位；
\item \path{scripts/ade_phase.py}：\texttt{phase\_from\_exact\_quaternions} 是 E 型入口，默认边长路线；
\texttt{phase\_details} 返回全部计算项，\texttt{method='angle'} 保留独立二面角比较；
\item \path{scripts/ade_count.py}：独立的 E 型精确计数代表，返回模群阶的整数指数；
与原几何接口的关系是式 \eqref{eq:count-conversion}，不是逐点相等；
\texttt{count\_details} 返回计数及指数，\texttt{algebraic\_phase} 返回复近似，
\texttt{gauge\_details} 返回转换平均的全部项；
\item \path{scripts/check_e_algebraic.py}：持久化精确整数五边形及完整余边界比较；
\item \path{scripts/check_ade_phase.py}：逐项五边形、手算例及退化例报告；
\item \path{docs/review-0922/SU2_HUMAN_REVIEW_RESPONSE.md}：本轮完整审计及可复制公式。
\end{itemize}
本轮还修复 $a^rb$ 坐标的末分量符号（应为 $-\sin(\pi r/n)$）、
E 入口静默截断非整数 $k$、A 包装函数在默认精度取倒数造成的精度丢失，
并声明 SymPy 验证依赖，修正 $\partial T$ 末项漏参数。
复现命令（先安装 \texttt{.[validation]}）：
\begin{verbatim}
python -m unittest discover -s tests -v
python -m scripts.su2_edge_crosscheck
python -m scripts.check_ade_phase
python -m scripts.check_e_algebraic
make adepdf
make completepdf
\end{verbatim}
当前 checkout 无 09-21 命名资料或原源 A 附件；本轮沿用已有 A 负号约定。
已给出 E 型完整有限几何公式及精确有界代数计数公式，
并证明显式余边界转换和完整类阶；E 型专属短表尚未导出。
任意 $n$ 的 D 表逐分支证明也仍需补齐。这些边界不以数值通过代替。
\end{sloppypar}

\ifdefined\SUComplete\expandafter\endinput\fi
\begin{thebibliography}{9}
\bibitem{jia} Q. Jia et al., \emph{Anomaly of Continuous Symmetries from Topological Defect Network}, arXiv:2510.14722v2 (2025).
\bibitem{murakami} J. Murakami, \emph{Volume formulas for a spherical tetrahedron}, arXiv:1011.2584v4 (2011).
\bibitem{epa} N. Epa and N. Ganter, \emph{Platonic and alternating 2-groups}, arXiv:1605.09192v1, Theorem 1.1.
\end{thebibliography}
\end{document}
